\documentclass{article}

\usepackage[preprint]{cs224r_2026}
\usepackage[utf8]{inputenc}
\usepackage[T1]{fontenc}
\usepackage{hyperref}
\usepackage{url}
\usepackage{booktabs}
\usepackage{amsfonts}
\usepackage{nicefrac}
\usepackage{microtype}
\usepackage{xcolor}
\usepackage{amsmath}
\usepackage{fancyhdr}
\usepackage{graphicx}
\usepackage{enumitem}

\pagestyle{fancy}
\fancyhf{}
\fancyhead[C]{CS224R Milestone Report}
\renewcommand{\headrulewidth}{0.4pt}

\renewcommand{\maketitle}{
  \begin{center}
    {\bf \large Fragment Assembly as Goal-Reaching: A Reinforcement Learning Approach to Targeted Molecular Design}\\
    \vspace{0.1cm}
    {\bf Team Members:} Katie Liu, Megan Santhumayor, Asmani Yamin\\
    \vspace{0.1cm}
    {\bf Emails:} katiel25@stanford.edu, msanth@stanford.edu, ayamin@stanford.edu\\
    \vspace{0.2cm}
  \end{center}
}

\begin{document}

\maketitle

\section{Experiments}

\textbf{Setup.}
We implemented the full goal-conditioned fragment assembly pipeline.
The environment maintains a partial molecule as a graph, initialized from a randomly sampled seed fragment from a library of 200 BRICS fragments (minimum 5{,}000 occurrences in M3-20M).
At each step the agent selects a fragment and attachment site from a BRICS-compatibility-filtered action set ($\sim$150 candidates).
The goal $\mathbf{g} \in [0,1]^3$ is a normalized vector of (LogP, QED, TPSA) sampled from 300 drug-like parent molecules.
The terminal reward is $r_T = -\|\phi_{\mathrm{norm}}(m) - \mathbf{g}\|_2 - 0.05 \cdot n_{\mathrm{open}}$, where $\phi_{\mathrm{norm}}$ is computed on the capped molecule and $n_{\mathrm{open}}$ penalizes unclosed attachment sites.
Partial-molecule state is encoded as a 512-bit Morgan fingerprint concatenated with $\mathbf{g}$ (515-dim total).
We train for 10{,}000 episodes with A2C and Hindsight Experience Replay (HER, $k{=}4$), using a replay buffer of 2{,}000 episodes.
Validation uses 200 greedy rollouts on a fixed held-out set; the metric is mean $\ell_2$ distance to goal.

\textbf{Baseline: A2C+HER (Morgan fingerprint state).}
The actor encodes state once then scores each candidate action independently; the critic estimates $V(s)$.
Both use 3-layer MLPs with hidden dimension 256.
Advantages are normalized per batch; entropy coefficient $\alpha = 0.005$.

\textbf{Experiment 1: PPO.}
We replaced the A2C update with the PPO clipped surrogate objective ($\epsilon{=}0.2$), collecting 16 on-policy episodes per update and running $K{=}4$ gradient epochs per batch with online HER relabellings.
All other hyperparameters were held fixed.

\textbf{Experiment 2: GNN state encoder.}
We replaced the Morgan fingerprint with a 3-layer MPNN operating on raw MolGraph arrays.
Node features (30-dim) encode element identity, formal charge, aromaticity, dummy-atom flag, and BRICS attachment-label; edge features (4-dim) encode bond type.
Mean pooling over a 128-dim node embedding produces a graph-level encoding concatenated with $\mathbf{g}$, then projected by the same MLP heads.
The A2C+HER training loop is otherwise identical; MolGraphs are stored in the replay buffer alongside transitions so HER relabellings copy them for free.

\textbf{Initial Results.}

\begin{table}[h]
  \centering
  \small
  \begin{tabular}{lcccc}
    \toprule
    Method & Val best & at ep & Val final & Train best \\
    \midrule
    Random policy    & 0.415 & ---  & ---   & --- \\
    A2C+HER          & 0.246 & 7000 & 0.258 & 0.227 \\
    PPO              & 0.332 & 9500 & 0.332 & 0.281 \\
    A2C+HER+GNN      & \textbf{0.241} & 8500 & 0.282 & 0.219 \\
    \bottomrule
  \end{tabular}
  \caption{Mean $\ell_2$ distance to goal on the held-out validation set (200 episodes, seed 99991). Lower is better. Random policy provides the untrained baseline.}
  \label{tab:results}
\end{table}

A2C+HER reduces mean distance from 0.415 (random) to 0.246 on validation.
The GNN encoder achieves the best result (0.241), improving over the fingerprint baseline by faster early learning (0.294 at ep~1000 vs.\ 0.337 for A2C+HER).
PPO underperforms significantly: policy entropy collapses from 4.98 to below 0.5 nats by episode 1{,}400 and never recovers, with validation distance plateauing at 0.33.
The collapse is driven by $K{=}4$ gradient epochs per batch over-fitting each rollout, compounded by HER transitions that carry approximate old log-probabilities, producing ill-conditioned importance ratios.

\begin{figure}[h]
  \centering
  \includegraphics[width=\linewidth]{FILENAME}
  \caption{Training (left) and validation (right) mean $\ell_2$ distance to goal for A2C+HER, PPO, and A2C+HER+GNN. The dotted line is the random-policy baseline (0.415). The GNN encoder achieves the lowest validation distance (0.241); PPO entropy collapse prevents meaningful improvement beyond episode 2{,}500.}
  \label{fig:comparison}
\end{figure}

\section{Changes to Research Hypothesis or Objective}

\textbf{Original hypothesis.}
Intermediate reward estimation (critic-shaped shaping and partial property proxies) would be the primary contribution, with terminal-reward+HER as a fallback.
The proposed training algorithm was off-policy (SAC or TD3).

\textbf{Revised hypothesis.}
HER alone provides sufficient dense training signal to overcome reward sparsity in goal-conditioned fragment assembly, making explicit intermediate reward estimation unnecessary for baseline convergence.
Furthermore, graph-structured state representations capture task-relevant information that fixed-size Morgan fingerprints cannot---specifically, which attachment sites are open and how fragments connect topologically---yielding measurably better goal-reaching performance.
Intermediate reward shaping is hypothesized to improve further upon the HER baseline, but is no longer the primary mechanism by which we address sparsity.

\textbf{Justification.}
SAC and TD3 are designed for continuous action spaces; our action set is discrete and changes size each step (150 candidates filtered by BRICS compatibility), making the continuous-action off-policy framework a poor fit without significant architectural modification.
A2C with a per-step action scorer handles variable-size discrete sets directly.
On the intermediate reward question, we implemented Crippen atom-contribution estimates of LogP on partial molecules and found them to be unreliable for short assemblies (2--4 atoms): contributions from open dummy sites introduce systematic bias, and partial QED and TPSA have no meaningful definition before a complete ring system is present.
HER achieves the same goal---dense feedback from incomplete episodes---by reinterpreting the achieved outcome as the target, which is exact rather than estimated.
These findings do not rule out intermediate rewards for longer assembly horizons and will be revisited in the final project.

\section{Next Steps}

\begin{enumerate}[leftmargin=*]
    \item \textbf{Further Experimentation:} Implement Crippen atom-contribution estimates of LogP on partial molecules as an auxiliary shaping reward $r_t^{\mathrm{partial}}$ and ablate against terminal-only and HER; train a scalar-reward (non-goal-conditioned) baseline optimizing a fixed linear combination of properties and evaluate on the same 1{,}000-target held-out set.
    \item \textbf{Algorithm Refinement:} Fix PPO's incompatibility with HER by tracking exact old log-probabilities for relabelled transitions and increasing the entropy coefficient ($\alpha \geq 0.02$) to prevent premature collapse; explore SAC adapted to the variable-size discrete action set as originally proposed.
    \item \textbf{Evaluation \& Analysis:} Compute success rate (fraction of generated molecules within threshold $\epsilon \in \{0.1, 0.2\}$ per property), chemical diversity of candidates per target, and validity rate; compare all methods against CPRL on the same target set as specified in the proposal evaluation plan.
\end{enumerate}

\textbf{Potential Challenges:} Partial LogP estimates via Crippen contributions are biased by open dummy sites, which contribute atom weights but no bonding context; disentangling this bias from a useful shaping signal may require masking dummy-atom contributions or normalizing by assembly progress. The GNN's actor loss shows high variance in late training (oscillating between $-25$ and $+30$), suggesting that gradient clipping at norm 1.0 is insufficient for the more expressive output distribution; tighter clipping or a lower actor learning rate will need to be tuned carefully to avoid destabilizing the policy. Finally, adapting SAC to a discrete, variable-size action set requires either a Gumbel-softmax relaxation or a discrete SAC variant, both of which add implementation complexity that may be difficult to complete within the remaining timeline.

\end{document}
